%% kapa: exam=1 round=29 year=2018 period=2 form=A subject=accounting subject_ko=회계학 seq=29 number=69
%% source: Q-Net 감정평가사 기출문제 (cst00309) artlSeq=5202808 file=extracted/1008/5202808/2018년 제29회 감정평가사 1차시험 기출문제/2018년 제29회 감정평가사 1차시험 기출문제/2교시 A형.pdf
%% conversion: kordoc 4.24.51 -> tools/gen1.py (auto); review=proofread-sonnet; flags=none
\begin{question}{69}
(주)감평은 20x3년 초 건물을 ￦41,500에 취득(내용연수 10년, 잔존가치 ￦1,500, 정액법 상각)하여 사용하고 있으며, 20x5년 중 손상이 발생하여 20x5년 말 회수가능액은 ￦22,500으로 추정되었다. 20x6년 말 건물의 회수가능액은 ￦26,000인 것으로 추정되었다. 동 건물에 대해 원가모형을 적용하는 경우 다음 설명 중 옳지 않은 것은?
\begin{choices}
  \item 20x5년 말 손상을 인식하기 전의 건물의 장부금액은 ￦29,500이다.
  \item 20x5년 건물의 손상차손은 ￦7,000이다.
  \item 20x6년 건물의 감가상각비는 ￦3,000이다.
  \item 20x6년 말 손상이 회복된 이후 건물의 장부금액은 ￦25,500이다.
  \item 20x6년 건물의 손상차손환입액은 ￦6,500이다.
\end{choices}
\end{question}
\answer{5}
